% H. Anexo tecnico: verificacion del detector y protocolo de seleccion
\chapter{Anexo técnico: verificación y selección del modelo}
\label{anexo:verificacion}

Este anexo recoge el detalle de tres procedimientos cuyos resultados se exponen de forma
resumida en el capítulo~\ref{cap:resultados}. Se separan del cuerpo por su extensión y porque
describen \textit{cómo} se obtuvo la evidencia, no la evidencia misma.

\section{Validación del nodo de sensado}

Los resultados de esta sección proceden de la campaña de caracterización EXP-002, de cuatro
minutos de duración sobre el activo en estado nominal. Se trata de una campaña destinada a
verificar la cadena de adquisición y no a producir el conjunto de entrenamiento, y así se
consigna.

El nodo integra los tres sensores en un mismo ciclo y publica en dos canales con cadencias
independientes. La cadencia del canal de ráfaga, que es la crítica porque determina la validez
del eje de frecuencias, se sostuvo en \mbox{1024 ms} en la totalidad de las ráfagas
registradas, sin desviación medible respecto del valor nominal. El cálculo completo de las
características de los tres ejes y de la ventana acústica consumió entre \mbox{29 ms} y
\mbox{30 ms}, frente a un periodo de \mbox{30 s} entre ráfagas.

El cuadro~\ref{table:validacion-nodo} recoge la repetibilidad de las características sobre el
activo en régimen estable, que es la magnitud que determina si un cambio de estado resultará
detectable.

\begin{table}[ht!]
\caption{Repetibilidad de las características en estado nominal (campaña EXP-002).}
\label{table:validacion-nodo}
\centering
\begin{tabular}{|p{4.6cm}|p{2.6cm}|p{2.2cm}|p{2.4cm}|}
\hline
\textbf{Magnitud} & \textbf{Media} & \textbf{Desviación} & \textbf{Coef. variación} \\
\hline
Frecuencia fundamental & \mbox{49,77 Hz} & 0,08 & 0,15 \% \\
\hline
Amplitud en la fundamental & \mbox{0,0660 m/s$^2$} & 0,0049 & 7,4 \% \\
\hline
\end{tabular}
\end{table}

La frecuencia fundamental se identificó en seis de las siete ráfagas registradas, siempre como
tercer pico espectral por detrás de los dos modos de la componente tonal descrita en el
capítulo~\ref{cap:diseno}. Su valor, \mbox{49,77 Hz}, es coherente con el régimen de un
motor de inducción de dos polos alimentado a \mbox{50 Hz} una vez descontado el
deslizamiento.

La estabilidad obtenida en frecuencia, con un coeficiente de variación del 0,15 \%,
constituye por sí misma una capacidad de medida: permite seguir variaciones pequeñas del
régimen de giro, que es uno de los indicadores de carga del compresor. La repetibilidad en
amplitud, del 7,4 \%, sitúa el umbral de detección en un cambio del orden del 21 \% para un
criterio de tres desviaciones típicas.

El canal térmico registró un salto de \mbox{5,7 °C} entre la superficie del compresor y el
ambiente, magnitud que discrimina sin ambigüedad el activo en marcha del activo detenido.
\section{Resiliencia frente al fallo del bus}

El cuadro~\ref{table:resiliencia} recoge el comportamiento de los mecanismos de integridad
durante la campaña EXP-002. Los contadores de salud descritos en el
capítulo~\ref{cap:diseno} permiten cuantificarlo sin instrumentación añadida.

\begin{table}[ht!]
\caption{Mecanismos de integridad del dato durante la campaña EXP-002, de cuatro minutos.}
\label{table:resiliencia}
\centering
\begin{tabular}{|p{5.6cm}|p{2.4cm}|p{4.0cm}|}
\hline
\textbf{Contador} & \textbf{Incremento} & \textbf{Interpretación} \\
\hline
Reintentos de lectura & 71 & Mediana de 1 por ráfaga, máximo de 30 \\
\hline
Rechazos por continuidad & 17 & Muestras corruptas detectadas \\
\hline
Ráfagas descartadas & 2 & De 9 registradas en total \\
\hline
Tramas del canal lento descartadas & 1 & De 259 registradas \\
\hline
Ráfagas no publicadas por falta de enlace & 0 & Enlace estable \\
\hline
\end{tabular}
\end{table}

El resultado relevante no es la magnitud de los contadores sino la relación entre ellos. Con
71 reintentos y 17 rechazos por continuidad, únicamente dos ráfagas se descartaron por
completo: el mecanismo de reintento absorbió la práctica totalidad de los fallos aislados del
bus sin sacrificar la ráfaga entera. Ello confirma el razonamiento de diseño expuesto en el
capítulo~\ref{cap:diseno}, según el cual descartar la ráfaga completa ante cualquier fallo
habría reducido la supervivencia a una fracción reducida del total.

Procede señalar un resultado no previsto, de utilidad directa para la preparación del conjunto
de datos. La ráfaga que registró 30 reintentos presentó también la kurtosis más elevada de la
ventana, 7,35 frente a una mediana de 4,69. Los contadores de salud correlacionan por tanto
con la calidad efectiva de la medida, lo que permite emplearlos como criterio objetivo de
filtrado del conjunto de datos en lugar de recurrir a una inspección manual.

Debe consignarse asimismo que las tramas descartadas no se contabilizan como valores nulos
sino como ausencia de dato, y que el análisis debe filtrarlas antes de calcular cualquier
estadístico, según el criterio establecido en el capítulo~\ref{cap:diseno}.
\section{Protocolo de selección libre de sesgo de espionaje}

El procedimiento descrito hasta aquí adolece de un defecto que procede consignar de forma
explícita, porque afecta a la validez de las cifras y no a su presentación. \textbf{La
característica sobre la que se construyó la regla se seleccionó ordenando las candidatas por su
capacidad de separación medida sobre el conjunto con fallo}, es decir empleando el conjunto de
evaluación para adoptar una decisión de diseño. Se trata de sesgo de espionaje de datos, y
cuanto de él se derive resulta optimista.

Se estableció por ello un protocolo con separación estricta. Su fundamento es la naturaleza del
detector: al ser de una clase, en explotación solo dispone de datos del activo sano, de donde se
sigue qué información es legítimamente utilizable.

Resulta admisible toda decisión adoptada sobre el conjunto nominal —en explotación también se
dispondría de él—, el conocimiento previo del dominio y las perturbaciones sintéticas de ese
conjunto. Resulta inadmisible seleccionar u ordenar características por su comportamiento sobre
el conjunto con fallo, ajustar umbrales atendiendo a la tasa de detección, y repetir la
evaluación final una vez observada.

La partición se realiza \textbf{por episodios y de forma cronológica}: los dieciséis más
antiguos para desarrollo y los ocho más recientes para prueba. El criterio cronológico reproduce
la situación real —se ajusta con lo capturado y se despliega sobre lo que venga después—
mientras que una partición aleatoria reparte episodios contiguos entre ambos lados y filtra
información del futuro.

El cuadro~\ref{table:protocolo} recoge el resultado.

\begin{table}[ht!]
\caption{Resultado del protocolo con separación estricta entre desarrollo y prueba.}
\label{table:protocolo}
\centering
\begin{tabular}{|p{7.0cm}|r|}
\hline
\textbf{Magnitud} & \textbf{Valor} \\
\hline
Modelo seleccionado sin emplear el conjunto con fallo & LOF \\
\hline
Direcciones de fallo cubiertas & 5 de 5 \\
\hline
Falsos positivos en el conjunto de desarrollo & 6,3 \% \\
\hline
\textbf{Falsos positivos en los episodios de prueba} & \textbf{8,3 \%} \\
\hline
\textbf{Detección sobre el activo con fallo} & \textbf{100 \%} \\
\hline
\end{tabular}
\end{table}

Debe destacarse que \textbf{el protocolo selecciona un modelo distinto} del que se había elegido
mediante el procedimiento anterior. Esa discrepancia constituye la evidencia de que aquella
elección estaba contaminada: de haber sido correcta, ambos procedimientos habrían coincidido.

El procedimiento inicial adolecía de un segundo defecto, de naturaleza distinta. Entre los
criterios de selección se incluyó la viabilidad de embarcar cada modelo, y varios candidatos se
descartaron por ese motivo \textbf{sin haberlo medido}. El cuadro~\ref{table:coste-nodo} recoge
la medición.

\begin{table}[ht!]
\caption{Coste real de embarcar cada modelo en el microcontrolador.}
\label{table:coste-nodo}
\centering
\begin{tabular}{|p{5.0cm}|r|r|}
\hline
\textbf{Modelo} & \textbf{Memoria} & \textbf{Operaciones por ráfaga} \\
\hline
Regla sobre n.\textsuperscript{o} de picos & \mbox{8 B} & 3 \\
\hline
\textit{Elliptic Envelope} & \mbox{960 B} & 240 \\
\hline
\textit{One-Class SVM} & \mbox{2,3 KB} & 1110 \\
\hline
LOF & \mbox{32 KB} & 15\,150 \\
\hline
\textit{Isolation Forest} & \mbox{274 KB} & 4800 \\
\hline
\end{tabular}
\end{table}

La plataforma dispone de \mbox{512 KB} de memoria estática y \mbox{8 MB} de memoria
pseudoestática, y \mbox{15\,150} operaciones en coma flotante a \mbox{240 MHz} corresponden al
orden de \mbox{0,1 ms} frente a una ráfaga cada \mbox{30 s}. \textbf{Los cinco modelos son
embarcables con amplio margen}, de modo que la restricción empleada como criterio de descarte no
existía y orientaba la selección hacia los modelos más simples sin fundamento.
\section{Cobertura frente a modos de fallo no observados}

La evidencia disponible corresponde a \textbf{un solo modo de fallo}. Dado que el sistema se
concibe como detector de anomalías y no como clasificador —se ajusta únicamente sobre el estado
nominal y señala lo que se desvía— cabría suponer que detectaría también otros modos, pero esa
suposición no está demostrada y procede acotarla.

Para localizar los puntos ciegos se realizó un análisis de sensibilidad: se perturbó el conjunto
nominal en las direcciones que la literatura asocia a modos de fallo característicos y se
observó qué candidatos reaccionaban. \textbf{Los datos así generados son sintéticos y no
constituyen evidencia experimental}; el análisis localiza puntos ciegos del detector, no
establece su rendimiento. El cuadro~\ref{table:cobertura-modos} recoge el resultado.

\begin{table}[ht!]
\caption{Análisis de sensibilidad frente a direcciones de fallo no observadas. Datos sintéticos.}
\label{table:cobertura-modos}
\centering
\begin{tabular}{|p{3.6cm}|r|r|r|r|r|}
\hline
\textbf{Dirección} & \textbf{Regla} & \textbf{\textit{LOF}} & \textbf{\textit{Ell. Env.}} & \textbf{Bosque} & \textbf{Vect. sop.} \\
\hline
Holgura mecánica & 100 \% & 100 \% & 100 \% & 100 \% & 100 \% \\
\hline
Rodamiento incipiente & \textbf{0 \%} & 100 \% & 100 \% & 100 \% & 100 \% \\
\hline
Roce o fricción & \textbf{0 \%} & 100 \% & 100 \% & 100 \% & 100 \% \\
\hline
Desequilibrio de masa & \textbf{0 \%} & 65 \% & 96 \% & \textbf{8 \%} & 66 \% \\
\hline
Obstrucción de ventilación & \textbf{0 \%} & 74 \% & \textbf{11 \%} & \textbf{10 \%} & 78 \% \\
\hline
\end{tabular}
\end{table}

El resultado obliga a matizar la elección de modelo expuesta más arriba. \textbf{La regla sobre
el número de picos resulta ciega a cuatro de las cinco direcciones examinadas}: detecta lo que
añade componentes espectrales y no lo que altera la amplitud o la impulsividad de la señal. Su
ventaja en tasa de falsos positivos se obtuvo midiendo contra el único modo de fallo disponible,
de modo que constituye una forma de sobreajuste a ese modo.

De ello se sigue una precisión sobre las características adimensionales: son \textbf{ciegas por
construcción} a un fallo que solamente modifique el nivel, puesto que toda magnitud se expresa
como cociente respecto de la fundamental. Un desequilibrio de masa eleva la amplitud a la
frecuencia de giro sin añadir componentes, y ninguna lo registra.

La corrección adoptada consiste en normalizar las magnitudes con unidades por la mediana del
\textit{propio} activo en funcionamiento, en lugar de suprimirlas. La característica resultante
es adimensional en su forma y específica en su valor, con lo que conserva la independencia de
máquina —la mediana vale 1 en ambos activos por construcción— y recupera la sensibilidad a la
amplitud: la \textit{Elliptic Envelope} pasa a detectar el 96 \% de las ráfagas perturbadas en esa
dirección y el \textit{Local Outlier Factor} el 65 \%. Debe consignarse la contrapartida operativa: el nodo necesita conocer la mediana de
su propio activo, de modo que ha de aprenderla durante una fase de referencia y no puede emitir
veredicto sobre su primera ráfaga.

Confrontar ese resultado con la validación por episodios pone de manifiesto un compromiso que
los datos disponibles no permiten resolver, y que el cuadro~\ref{table:compromiso} resume.

\begin{table}[ht!]
\caption{Compromiso entre especificidad medida y cobertura de modos de fallo.}
\label{table:compromiso}
\centering
\begin{tabular}{|p{5.0cm}|r|r|}
\hline
\textbf{Modelo} & \textbf{Peor episodio} & \textbf{Direcciones cubiertas} \\
\hline
Regla sobre n.\textsuperscript{o} de picos & \textbf{2,7 \%} & 1 de 5 \\
\hline
\textit{Elliptic Envelope} & 20,0 \% & 4 de 5 \\
\hline
Envolvente de Mahalanobis & 20,0 \% & 4 de 5 \\
\hline
\textbf{\textit{Local Outlier Factor}} & \textbf{28,6 \%} & \textbf{5 de 5} \\
\hline
\textit{Isolation Forest} & 46,8 \% & 3 de 5 \\
\hline
\textit{One-Class SVM} & 71,4 \% & \textbf{5 de 5} \\
\hline
\end{tabular}
\end{table}

La regla es diez veces mejor en especificidad y cubre la quinta parte de las direcciones. No hay
en estos datos ningún candidato que gane en ambas columnas, y la elección entre ellas no es
técnica sino de política de operación: depende de si el coste de una alarma falsa supera al de un
fallo no detectado, que es una decisión del explotador del activo y no del diseñador del sistema.
El criterio declarado de antemano fue la máxima cobertura y, a igual cobertura, la menor tasa de
falsos positivos, lo que selecciona el \textit{Local Outlier Factor} frente a la
\textit{One-Class SVM}. Las cifras de esta tabla proceden del conjunto de desarrollo, dado que la
selección debía resolverse sin mirar el de prueba.
\section{Implementación de la inferencia sin entorno de ejecución}

De lo anterior se sigue una consecuencia sobre la implementación en el nodo que conviene
enunciar, dado que contradice la expectativa habitual en este tipo de trabajos. Los entornos de
ejecución de modelos reducidos para microcontroladores, según se expuso en el
capítulo~\ref{cap:estado-arte}, operan sobre \textbf{redes neuronales}. Ninguno de los
candidatos considerados pertenece a esa familia: el método seleccionado se basa en vecindad, la
envolvente es una forma cuadrática y el \textit{Isolation Forest} un conjunto de árboles de
decisión. No se trata por tanto de que un entorno de inferencia resulte excesivo para el modelo
adoptado, sino de que \textbf{no puede ejecutarlo}.

Debe añadirse que una red neuronal no constituía una alternativa viable: con 505 observaciones
nominales de veinticuatro episodios, cualquier arquitectura con capacidad apreciable se
sobreajustaría. La familia de modelos quedó determinada por el tamaño de muestra y no por la
plataforma.

La inferencia se programa en consecuencia directamente en el nodo —del orden de quince líneas
para la envolvente y sesenta para el método seleccionado— y los parámetros se generan desde el
entorno de análisis en un fichero de cabecera, con anotación de la campaña y la fecha de las que
proceden, de modo que reentrenar consista en volver a ejecutar el procedimiento y recompilar.

La ejecución del modelo en el microcontrolador se implementó en C++ y se verificó en dos
niveles complementarios: primero frente a la implementación de referencia en el entorno de
análisis, y posteriormente mediante despliegue y ejecución continua sobre la plataforma física
ESP32-S3 en la campaña EXP-007 (\mbox{21,54 h} de operación ininterrumpida).

Las rutinas del detector se escribieron \textbf{sin dependencias del entorno de ejecución del
microcontrolador}, siguiendo el criterio ya aplicado a la extracción de características: ello
permite compilarlas en un computador y contrastarlas frente a la implementación de referencia
antes de emplear un ciclo de programación de la placa. Los parámetros ajustados se generan en un
fichero de cabecera independiente de las rutinas, de modo que reentrenar el modelo no obligue a
modificar código.

La verificación cruzada previa sobre las \textbf{1161 ráfagas reales} de las dos primeras
campañas arrojó coincidencia plena (0 discrepancias sobre 1161 casos, con un error relativo de
\mbox{2,1$\cdot$10$^{-7}$} en LOF), tal como resume el cuadro~\ref{table:verificacion-nodo}.

\begin{table}[ht!]
\caption{Verificación del código embarcado frente a la implementación de referencia.}
\label{table:verificacion-nodo}
\centering
\begin{tabular}{|p{5.6cm}|r|r|}
\hline
\textbf{Magnitud} & \textbf{Modelo principal} & \textbf{Envolvente} \\
\hline
Error relativo en la puntuación & \mbox{2,1$\cdot$10$^{-7}$} & \mbox{3,6$\cdot$10$^{-4}$} \\
\hline
Margen del caso más próximo al umbral & --- & 9,2 veces el error \\
\hline
\textbf{Veredictos discrepantes} & \textbf{0 de 1161} & \textbf{0 de 1161} \\
\hline
\end{tabular}
\end{table}

Tras flashear el firmware en el nodo físico y desplegar la captura continua durante
\mbox{21,54 h} (campaña EXP-007), se contrastaron las 911 ráfagas evaluables medidas en la placa
con el recálculo idéntico en el servidor. El resultado ratifica la paridad: \textbf{906 de 911
veredictos coincidentes (99,45~\%)}. Las cinco discrepancias (0,55~\%) corresponden a
observaciones situadas en la frontera inmediata del umbral de decisión, donde pequeñas
divergencias de precisión en coma flotante determinan el corte. La afirmación de que el
diagnóstico se ejecuta en el nodo queda por tanto \textbf{medida experimentalmente en hardware}
y no supuesta.

Procede extender el contraste a la operación completa, dado que el tramo perturbado lo degrada.
Sobre las \mbox{29,09 h} y sus 1340 ráfagas evaluables la paridad baja al \mbox{98,81 \%} —1324
coincidentes, 16 discrepancias— y el error máximo en la puntuación sube de \mbox{0,132} a
\mbox{0,880}. Debe declararse la consecuencia: ese error máximo es del orden del margen que
separa del umbral a la observación más próxima, de modo que la coincidencia entre las dos
implementaciones \textbf{no está garantizada para observaciones limítrofes}. La causa no es un
defecto de la derivación de características —anular el gradiente térmico empeora la coincidencia
en lugar de mejorarla— sino la precisión en coma flotante de treinta y dos bits del
microcontrolador frente a la de sesenta y cuatro del equipo de análisis.

El cuadro~\ref{table:coste-inferencia} recoge el coste medido de cada alternativa.

\begin{table}[ht!]
\caption{Coste de embarcar cada modelo, medido sobre los parámetros ajustados.}
\label{table:coste-inferencia}
\centering
\begin{tabular}{|p{5.0cm}|r|r|}
\hline
\textbf{Modelo} & \textbf{Memoria} & \textbf{Operaciones} \\
\hline
Regla sobre n.\textsuperscript{o} de picos & \mbox{8 B} & 3 \\
\hline
\textit{Elliptic Envelope} & \mbox{840 B} & 210 \\
\hline
\textit{One-Class SVM} & \mbox{2,3 KB} & 1110 \\
\hline
Modelo seleccionado (LOF) & \mbox{31,6 KB} & 15\,150 \\
\hline
\textit{Isolation Forest} & \mbox{274 KB} & 4800 \\
\hline
\end{tabular}
\end{table}

La medición empírica del tiempo de ejecución en el microcontrolador ESP32-S3 a \mbox{240 MHz}
(registrada campo a campo mediante el contador \texttt{us\_inference}) confirmó con holgura estas
cifras. Para las 2536 ejecuciones de la campaña EXP-007, el tiempo de inferencia arrojó una
\textbf{mediana de \mbox{14 \textmu s}}, un percentil 99 de \mbox{1380 \textmu s} (1,38~ms) y un
máximo absoluto de \mbox{1942 \textmu s} (1,94~ms). Frente al periodo de muestreo de una ráfaga
cada \mbox{30 s}, la inferencia consume únicamente el \textbf{0,00005~\% del ciclo de la CPU},
demostrando que el coste computacional del detector no supone limitación alguna en el borde.

\section{Evaluación de la política de avisos}

reinicia su cuenta ante una única ráfaga nominal, y durante la anomalía sostenida se intercalaron
tres, con lo que el sistema volvió a notificar en cada ocasión. En operación son cuatro
notificaciones de lo mismo.

Procede señalar aquí una propiedad del diseño que resulta ventajosa. La histéresis es una
\textbf{política de notificación y no una medida}: no altera el veredicto ni las puntuaciones,
solo el aviso y la cuenta de racha. Y al ser una función pura de la secuencia de veredictos,
cualquier política alternativa se evalúa \textbf{sobre los datos ya capturados}, sin modificar el
firmware ni repetir ninguna campaña. El cuadro~\ref{table:politica} recoge esa evaluación.

\begin{table}[ht!]
\caption{Evaluación de un tiempo mínimo entre avisos sobre los veredictos capturados.}
\label{table:politica}
\centering
\begin{tabular}{|r|r|r|l|}
\hline
\textbf{Silencio} & \textbf{Avisos} & \textbf{Del episodio} & \textbf{1.\textsuperscript{a} detección} \\
\hline
0 min (actual) & 18 & 4 & 12:10:40 \\
\hline
30 min & 12 & 3 & 12:10:40 \\
\hline
60 min & 12 & 2 & 12:10:40 \\
\hline
\textbf{120 min} & \textbf{7} & \textbf{1} & 12:10:40 \\
\hline
\end{tabular}
\end{table}

Debe consignarse que se descartó una primera propuesta: exigir varias ráfagas
\textit{nominales} consecutivas para rearmar no reduce los avisos —14, 16 y 19 según el valor,
frente a los 18 de la política actual— porque al no romper la racha con una nominal aislada, la
de anomalías se acumula con mayor rapidez y el aviso se dispara antes.

Lo que sí resulta eficaz es el tiempo mínimo entre avisos, que constituye deduplicación y no
insensibilización: la primera detección de cada episodio no se pierde con ningún valor ensayado.
La contrapartida ha de declararse, dado que un fallo aparecido dentro de la ventana de silencio
quedaría sin notificar hasta que esta expirase, y con un ciclo de trabajo de unos \mbox{85 min}
un silencio de \mbox{120 min} abarca más de un ciclo completo. Es por tanto una decisión de
política de operación y no una elección técnica.

No se modificó el firmware: alterar la política a mitad de una serie haría que el campo de aviso
significase algo distinto antes y después del cambio.

Con independencia de lo anterior, dos resultados de esta campaña son sólidos y ambos afectan al diseño. El primero es que
\textbf{el mecanismo de histéresis se comporta según su diseño}: se acumularon exactamente tres
ráfagas anómalas consecutivas y el aviso se emitió en la tercera y no antes, según se aprecia en
la figura~\ref{fig:puerta}. El segundo es la cuantificación del coste del canal acústico. Anulando las tres bandas y
recalculando la puntuación sobre las ráfagas de la captura cerrada a las \mbox{12:58} del 28 de
agosto, la tasa de ráfagas
marcadas desciende del \textbf{14,0 \% al 5,7 \%}, y el episodio de la apertura deja de
detectarse por completo: ninguna de sus tres ráfagas resulta marcada. El canal acústico aporta
en consecuencia \textbf{8,3 puntos} de tasa de ráfagas marcadas y es el único responsable de la
detección de la apertura.

Debe matizarse el mecanismo: únicamente el 10 \% de las ráfagas anómalas presenta un valor
extremo en la banda intermedia, frente al 5 \% que correspondería por construcción, de modo que
la contribución no procede de sucesos intensos como la apertura sino de una dispersión que la
campaña de referencia no capturó.

El canal acústico presenta por tanto un \textbf{compromiso con las dos caras medidas}: aportó la
confirmación independiente del fallo real, procedente de un sensor que no comparte cadena con el
acelerómetro, y cuesta 8,3 puntos de tasa de marcado frente a perturbaciones del entorno. La
elección entre ambas depende de si el activo opera en un entorno acústicamente estable, de modo
que queda declarada y no resuelta. Los datos sugieren una alternativa sin ensayar: exigir que la
desviación se manifieste \textit{también} en el canal de vibración antes de notificar.
